% BHU PYQ -- Professional Exam Template
\documentclass[11pt,a4paper]{article}

\usepackage[a4paper,top=2.5cm,bottom=2.5cm,left=2.8cm,right=2.8cm,headheight=14pt]{geometry}
\usepackage{amsmath,amssymb}
\usepackage[T1]{fontenc}
\usepackage[utf8]{inputenc}
\usepackage{lmodern}
\usepackage{microtype}
\usepackage{fancyhdr}
\usepackage{enumitem}
\usepackage{hyperref}
\usepackage{lastpage}
\usepackage{parskip}

\hypersetup{colorlinks=true,urlcolor=black,linkcolor=black,
  pdftitle={B.Sc. (Hons.) Botany Sem-V BOB-503 -- BHU PYQ}}

\pagestyle{fancy}
\fancyhf{}
\renewcommand{\headrulewidth}{0.4pt}
\renewcommand{\footrulewidth}{0.4pt}
\fancyhead[L]{\small   \textit{Banaras Hindu University}}
\fancyhead[R]{\small   \textit{Botany Paper: BOB-503}}
\fancyfoot[C]{\small Page  \thepage\ of \pageref{LastPage}}


\newlist{parts}{enumerate}{2}
\setlist[parts,1]{label=(\alph*),leftmargin=*,itemsep=4pt,topsep=3pt}
\setlist[parts,2]{label=(\roman*),leftmargin=*,itemsep=2pt,topsep=2pt}

\newcommand{\pts}[1]{\hfill   \textit{[#1]}}


\begin{document}


\begin{center}
  {\large\bfseries BANARAS HINDU UNIVERSITY}\\[2pt]
  {\normalsize B.Sc. (Hons.) Semester V Examination, 2016-17}\\[6pt]
  \rule{0.8\linewidth}{0.5pt}\\[6pt]
  {\large\bfseries Botany}\\[4pt]
  {\large\bfseries Paper: BOB-503 --- Plant Ecology}\\[4pt]
  {\normalsize Time: Three Hours \hfill Full Marks: 70}
\end{center}

\rule{\linewidth}{0.4pt}

\smallskip



\noindent   \textit{Note: Attempt five questions in all, selecting at least two questions from Section-A and Section-B. Question No. 1 is compulsory. The figures in the right-hand margin indicate marks.}

\bigskip

\medskip


\noindent   \textbf{Question 1.} Answer the following in not more than 50 words each:\pts{$1  \times 10 = 10$}

\begin{enumerate}[label=(\roman*),leftmargin=*,itemsep=4pt,topsep=3pt]
  \item Define ecosystem.
  \item Name two major greenhouse gases.
  \item Stratification in water body.
  \item What is soil profile?
  \item Define carrying capacity.
  \item Decibel is a unit of what?
  \item Minamata disease is caused by \makebox[3cm]{\hrulefill}.
  \item Photochemical smog.
  \item Give two examples of primary succession.
  \item Humus.
\end{enumerate}

\bigskip

\begin{center}
   \textbf{SECTION -- A}
\end{center}

\medskip


\noindent   \textbf{Question 2.} What is a plant community? Discuss synthetic characters of a community.\pts{15}

\medskip


\noindent   \textbf{Question 3.} Define population. Discuss various characters of population.\pts{15}

\medskip


\noindent   \textbf{Question 4.} Describe the nitrogen cycle in detail.\pts{15}

\medskip


\noindent   \textbf{Question 5.} Write critical notes on any two of the following:\pts{$7.5  \times 2 = 15$}

\begin{parts}
  \item Light as ecological factor
  \item Carbon cycle
  \item Universal model of energy flow
\end{parts}

\bigskip

\begin{center}
   \textbf{SECTION -- B}
\end{center}

\medskip


\noindent   \textbf{Question 6.} Write critical comments on any two of the following:\pts{$7.5  \times 2 = 15$}

\begin{parts}
  \item Eutrophication
  \item Importance Value Index (IVI)
  \item Concept of ecological niche
\end{parts}

\medskip


\noindent   \textbf{Question 7.} Define ecological community. Discuss various theories of mechanism of ecological succession.\pts{15}

\medskip


\noindent   \textbf{Question 8.} Write descriptive notes on any two of the following:\pts{$7.5  \times 2 = 15$}

\begin{parts}
  \item Acid rain
  \item Thermal pollution
  \item Teratogens
\end{parts}

\medskip


\noindent   \textbf{Question 9.} Define biodiversity. Discuss in brief indices of species diversity.\pts{15}

\vfill
\rule{\linewidth}{0.4pt}

\begin{center}\small
\href{https://bhu-exam.vercel.app/}{Click Here}\\[4pt]\href{https://me-aryan.vercel.app/}{Click Here}\\[4pt]\href{https://quashan-me.vercel.app/}{Click Here}
\end{center}

\end{document}
